% cap2.tex

\chapter{Marco Teórico}
\label{c2}

% ============================================================
% Contenido retomado de Proyecto de Título I (PT1), Hito 1 (Hito_1.tex)
% ============================================================

El presente capítulo contiene el marco teórico necesario para sustentar el análisis de redes financieras. De esta forma, se estudiará, bajo el marco teórico presentado, la competencia que surge entre infraestructuras tradicionales y alternativas.

\section{Modelos de redes financieras}\label{c22}

Se definen las redes como conexiones entre nodos y \textit{links}: los agentes
(nodos) toman decisiones estratégicas sobre con quiénes formar vínculos,
considerando los beneficios esperados y sus costos \cite{jackson_strategic_1996}. Con base en estas redes se modelan fenómenos como el contagio financiero, la propagación del riesgo sistémico y la estabilidad del sistema económico. La literatura distingue tres
modelos de referencia:

\begin{itemize}
    \item \textbf{\citeA{eisenberg_systemic_2001}:} Representan las redes
    financieras a través de una matriz de obligaciones entre agentes. El modelo se distingue pues captura el efecto
    cascada en la red cuando ocurre un \textit{default}, mostrando la existencia
    y unicidad de un vector de pagos que liquida el sistema.

    \item \textbf{\citeA{elliott_financial_2014}:} Modelan el contagio financiero en dos dimensiones. La primera dimensión corresponde a la dependencia de contraparte (integración) y la segunda por el número de contrapartes por agente (diversificación). Se distingue en el modelo que ambas dimensiones tienen efectos no monótonos sobre la amplificación o mitigación del riesgo sistémico.

    \item \textbf{\citeA{jalan_incentive-aware_2024}:} Introducen una matriz cuadrada para representar origen-destino junto con otra matriz $Q$, simétrica y semidefinida positiva, que captura la estructura de costos y
    beneficios entre pares de agentes. Además utiliza una función de utilidad media-varianza que permite modelar decisiones óptimas de contratación bilateral en un equilibrio \textit{pairwise stable}. Es el modelo base sobre
    el cual se construye la tesis de \citeA{venegas_reyes_financial_2025}, que este
    trabajo extiende.
\end{itemize}

\paragraph{Propagación de riesgo y estabilidad}

Las redes financieras tradicionales son particularmente vulnerables al contagio
cuando falla un nodo crítico, como un banco sistémico. Como muestra \citeA{glasserman_contagion_2016}, las redes densamente conectadas pueden ser resilientes ante \textit{shocks} pequeños pero frágiles ante crisis sistémicas. Demuestra que la conectividad absorbe perturbaciones locales, pero también amplifica los choques de gran magnitud. En la misma línea, \citeA{acemoglu_systemic_2015} documentan que la diversificación de conexiones
incrementa inicialmente la estabilidad del sistema, pero más allá de cierto
umbral la interdependencia generada abre nuevos canales de riesgo sistémico. Es
decir, la relación entre conectividad y estabilidad no es monótona.

Las redes financieras tradicionales presentan, además, limitaciones estructurales
documentadas extensamente en la literatura:

\paragraph{Altos costos de intermediación}
Los intermediarios financieros tales como bancos o cámaras de compensación imponen comisiones por transferencias internacionales y por gestión de riesgo y cumplimiento regulatorio que en promedio alcanzan un $6,36\%$ del monto enviado a nivel mundial (más del doble de la meta de $3\%$ fijada en el ODS 10.c), cifra que baja a $4,59\%$ cuando la transferencia se realiza por canales digitales frente a $7,30\%$ en canales no digitales \cite{world_bank_rpw_2025}. Esto constituye una barrera a la inclusión financiera, en particular en mercados emergentes.

\paragraph{Opacidad y asimetrías de información}
La información transaccional suele restringirse a las partes directamente
involucradas, lo que dificulta evaluar el riesgo de una contraparte y favorece
el comportamiento oportunista.

\paragraph{Riesgo sistémico y centralización}
La centralización en instituciones clave como grandes bancos y mercados
centralizados crea puntos únicos de falla. \citeA{elliott_financial_2014}
señalan que esta concentración del riesgo hace a las redes tradicionales más
susceptibles de propagar shocks macroeconómicos o eventos de liquidez. En
contraste, \citeA{gofman_efficiency_2017} argumentan que los bancos grandes e
interconectados pueden mejorar la eficiencia agregada. Esto, precisamente porque
reducen el largo de las cadenas de intermediación en el mercado. El efecto neto depende de si domina la reducción de fricciones de intermediación o la concentración del riesgo, de modo que la centralización no es unívocamente negativa.

\paragraph{Baja inclusión financiera}
Las redes tradicionales suelen excluir a agentes que no cumplen criterios
formales de acceso como un historial crediticio, cuenta bancaria. Un problema
particularmente agudo en economías en desarrollo, donde aproximadamente 2 mil millones de adultos (un $38\%$ de la población adulta mundial, cifra de 2014) permanecían fuera del sistema financiero formal \cite{singer_financial_2017}.

\paragraph{Efectos de red}
Una mayor conectividad amplía el acceso a financiamiento y permite diversificar
riesgo. Por otra parte, también puede generar congestión operacional y mayores costos transaccionales a medida que la red crece \cite{hinzen_bitcoins_2022,iyengar_economics_2023}.

\section{Competencia entre infraestructuras financieras}\label{c23}

La competencia entre redes Tradicional y Alternativa que este trabajo analiza, que \citeA{venegas_reyes_financial_2025} modela formalmente, se apoya en dos
cuerpos de literatura complementarios. El primero, la teoría de plataformas y mercados de
dos lados que explica cómo compiten infraestructuras con efectos de red. Segundo, el arbitraje regulatorio que explica por qué surgen infraestructuras alternativas.

\subsection{Mercados de dos lados y competencia entre plataformas}

\citeA{economides_economics_1996} establece que las redes financieras
funcionan económicamente como plataformas debido a externalidades de red. La
utilidad de una conexión financiera aumenta con la liquidez y profundidad de la
red a la que pertenece. \citeA{rochet_platform_2003} formalizan esta idea en el
marco de mercados de dos lados. Este concepto se define como mercados donde una plataforma
habilita interacciones entre dos grupos de usuarios y donde el valor para un
lado depende críticamente de la participación del otro lado (externalidades de
red \textit{indirectas}). Traducido en el contexto de redes financieras, estos dos lados se pueden entender como oferentes de liquidez (capital disponible para prestar o
invertir) y demandantes de liquidez (agentes que buscan financiamiento o
cobertura).

Por su parte, \citeA{armstrong_competition_2006} extiende este marco a la competencia entre
múltiples plataformas e identifica tres escenarios según el patrón de
participación de los agentes. \textit{single-homing}, los usuarios se afilian a
una sola plataforma, lo que tiende a generar resultados de tipo
\textit{winner-take-all}. \textit{multi-homing}, donde los usuarios se afilian a varias plataformas simultáneamente, lo que suaviza la competencia y permite la coexistencia. Por último, en los \textit{cuellos de botella
competitivos}, un lado se afilia a una sola plataforma mientras el otro se
multi-afilia. El resultado central de Armstrong es que la coexistencia de
infraestructuras es viable cuando el costo de multi-homing es moderado y las
plataformas se diferencian. Ya que, si el costo es muy alto, el mercado tiende a concentrarse en una sola infraestructura.

Este marco aplica directamente al modelo dual de \citeA{venegas_reyes_financial_2025}: la variable de decisión $(w_i^T, w_i^A)$ es, finalmente, una decisión de multi-homing entre la red Tradicional y la Alternativa. El parámetro de
fricción $\lambda_A$ representa el costo de ese multi-homing. \citeA{caillaud_chicken_2003} muestran que la diferenciación entre plataformas actúa como un mecanismo que suaviza la competencia directa. Esta diferenciación puede darse por ejemplo en la transparencia o en el tipo de contratos que cada red permite. De esta forma se sostiene la coexistencia entre plataformas, en lugar de forzar un resultado de dominancia total de una infraestructura sobre la otra.

\subsection{Arbitraje regulatorio}

Un elemento central para entender por qué surgen infraestructuras alternativas es el concepto de \textbf{arbitraje regulatorio}. Se conceptualiza esta dinámica por \citeA{kane_good_1977} como la ``dialética regulatoria''. Quien a su vez la define como un proceso cíclico en que la regulación induce a los agentes económicos a innovar tecnológicamente, con el único fin de eludir restricciones costosas. A diferencia de modelos estáticos donde los agentes aceptan pasivamente las restricciones, el volumen ``excluido'' busca nuevos canales, dando origen a un sistema ``en la sombra''.

Esta dinámica está documentada empíricamente por \citeA{buchak_fintech_2018}, quienes encuentran evidencia causal de que el arbitraje regulatorio es un motor determinante del crecimiento \textit{Fintech}. 
Estiman que cerca del $50\%$ de la expansión del \textit{shadow banking} responde a la carga regulatoria sobre la banca tradicional (mayores requerimientos de capital, entre otros).
Por lo tanto, no se trata solo de una decisión de preferencia por la tecnología que la acompaña.
Esto valida el mecanismo central del modelo dual, que consiste en que la migración de volumen hacia la red Alternativa ($w_i^A$) es una respuesta económica racional a la fricción regulatoria en el sector incumbente, no meramente una preferencia tecnológica.
La literatura, sin embargo, no es unánime sobre las implicancias de bienestar de
este arbitraje. \citeA{buchak_fintech_2018} lo interpretan como una
respuesta eficiente a restricciones costosas. Por su parte, \citeA{gorton_regulating_2010}
advierten que el \textit{shadow banking} puede ser también una fuente de
inestabilidad sistémica, precisamente por operar con menor transparencia y
supervisión que el sistema regulado. El parámetro de transparencia $\tau$ en el
modelo dual puede leerse como una forma de capturar esta tensión. Una infraestructura alternativa transparente ($\tau > 0$) podría ofrecer los beneficios de eficiencia del arbitraje regulatorio, sin heredar necesariamente
los riesgos típicamente asociados al \textit{shadow banking} opaco.

\subsection{Adopción tecnológica, efectos de red y dependencia de trayectoria}

La adopción de una infraestructura financiera alternativa enfrenta, además,
problemas de coordinación característicos de mercados con externalidades de
red. \citeA{katz_network_1985} muestran que las tecnologías con externalidades
de red positivas, donde el valor para cada usuario aumenta con el número de
otros usuarios, enfrentan un problema de coordinación fundamental. Los adoptantes potenciales pueden preferir no adoptar si esperan que pocos otros lo hagan, incluso cuando la adopción universal beneficiaría a todos. 
Esto genera \textit{exceso de inercia}, que consiste en que el mercado permanece anclado a una tecnología
inferior porque los agentes esperan racionalmente que otros sigan usando el
estándar establecido. También causa equilibrios múltiples para los mismos parámetros
fundamentales, dependiendo únicamente de las expectativas sobre el
comportamiento de los demás.

\citeA{farrell_installed_1986} formalizan el concepto de \textit{exceso de
inercia} sobre cómo una tecnología establecida con una base
de usuarios grande puede resistir la competencia de alternativas objetivamente
superiores, porque cada agente racionalmente espera que otros se muevan
primero. La red Tradicional posee, en este sentido, una base ya instalada con relaciones existentes, infraestructura, marcos regulatorios y procedimientos operacionales, base que la red Alternativa debe superar con beneficios de transparencia ($\tau$) suficientemente grandes para compensar tanto los costos de fricción directos ($\lambda_A$) como la ventaja de coordinación de la red incumbente.

\citeA{arthur_competing_1989} extiende esta lógica introduciendo los conceptos
de \textit{path dependence} y \textit{lock-in} en mercados con retornos
crecientes a la adopción. Cuando una tecnología exhibe retroalimentaciones
positivas, pequeños accidentes históricos o ventajas tempranas pueden llevar a
la dominancia permanente de una tecnología, incluso si otra es objetivamente
superior. Aplicado al modelo dual, esto implica que la asignación inicial
$w_i^*$, heredada del equilibrio de red única sin restricciones, condiciona el
punto de partida desde el cual se evalúa la adopción de la red Alternativa. La
multiplicidad de equilibrios del modelo dual (cuando $Z_F^{\text{dual}}$ no
tiene rango completo) puede interpretarse precisamente como una manifestación de
esta dependencia de trayectoria.

\section{Blockchain y finanzas descentralizadas}

La tecnología \textit{blockchain}, introducida por \citeA{nakamoto_bitcoin_2008}
como la estructura subyacente de Bitcoin, es un sistema de registro
distribuido que permite consignar transacciones de forma segura, transparente e
inmutable. Esto se logra a través de una red de nodos, sin requerir una autoridad central que
controle el libro de cuentas. \citeA{cong_blockchain_2019} identifican cuatro
mecanismos económicos centrales a través de los cuales blockchain afecta los
resultados de mercado:

\begin{itemize}
    \item \textbf{Transparencia y auditabilidad}, ya que todas las transacciones quedan registradas en un libro público e inmutable, verificable por cualquier participante.
    \item \textbf{Automatización vía \textit{smart contracts}}, que permite ejecutar términos contractuales sin depender de intermediarios de confianza, reduciendo el riesgo de contraparte.
    \item \textbf{Desintermediación}, al habilitar transacciones directas entre pares.
    \item \textbf{Composabilidad}, que permite construir productos financieros complejos combinando contratos más simples.
\end{itemize}

Estos mismos autores advierten que estas propiedades introducen nuevas disyuntivas entre eficiencia y otros
objetivos, como privacidad o flexibilidad contractual que no están presentes en
las redes tradicionales.

La evidencia empírica sobre blockchain como infraestructura alternativa es
mixta. \citeA{makarov_blockchain_2021}, en un análisis extenso del mercado de
Bitcoin, documentan que la transparencia transaccional reduce ciertas asimetrías
de información. Los participantes sofisticados pueden observar patrones de flujo
y comportamiento de las contrapartes en tiempo real. Simultáneamente crea
nuevas ventajas de información para quienes tienen mayor capacidad analítica
sobre los datos \textit{on-chain}. Además, encuentran que los costos de
transacción son altamente variables y pueden dispararse durante periodos de
congestión de red. \citeA{huberman_monopoly_2021} caracterizan este fenómeno
como un ``monopolio sin monopolista''. El protocolo restringe la oferta de
espacio de bloque, generando precios de monopolio, pero ninguna entidad
individual captura esas rentas (que van, en cambio, a los mineros o validadores).
Esta evidencia empírica sobre costos variables y crecientes con la congestión es
consistente con la forma funcional cuadrática $\lambda_A \|w_i^A\|^2$ que el
modelo dual utiliza para representar el costo de fricción de la red Alternativa.

\citeA{schar_decentralized_2021} identifica los canales concretos a través de
los cuales la transparencia de blockchain crea valor económico. En primer lugar, reduce la
selección adversa en mercados \textit{over-the-counter} al hacer observables
los historiales transaccionales. Además, mejora el monitoreo de riesgo de contraparte.
También habilita la ejecución \textit{trustless} de operaciones complejas sin depender
de intermediarios humanos. Por último, facilita la auditoría regulatoria. Sin embargo,
\citeA{werbach_trust_2018} matiza esta lectura. La evidencia de vulnerabilidades
en \textit{smart contracts}, la persistencia de mecanismos de gobernanza humana
para resolver disputas, y el llamado ``problema del oráculo'' (la necesidad de
información externa confiable para que los contratos inteligentes funcionen)
muestran que blockchain no elimina la confianza. En realidad, la redistribuye de
la confianza en instituciones humanas a la confianza en código, criptografía y
mecanismos de consenso. Sus beneficios dependen del contexto institucional
específico. Estos mecanismos, transparencia con beneficios acotados por
limitaciones prácticas, son los que el parámetro $\tau$ del modelo dual busca
capturar de forma neta.

Finalmente, la diferencia regulatoria entre redes Tradicional y Alternativa
($\Psi_i^T \neq \Psi_i^A$) tiene respaldo institucional directo.
\citeA{zetzsche_decentralized_2020} documentan que los sistemas
descentralizados suelen operar en zonas grises regulatorias. Esto surge por la dificultad
de determinar qué jurisdicción aplica a un protocolo global, lo que habilita cierto arbitraje jurisdiccional pero también genera incertidumbre legal. 
\citeA{auer_rise_2023}, en un estudio del Banco de Pagos Internacionales, documentan un patrón de adopción institucional \textit{bifurcado}. 
Las instituciones financieras grandes y reguladas experimentan principalmente con
\textit{blockchains} permisionados que preservan privacidad y cumplimiento
regulatorio. 
Mientras que entidades más pequeñas o menos reguladas adoptan más directamente sistemas públicos y sin permisos. 
En ambos casos, la adopción se concentra en casos de uso específicos como pagos transfronterizos, financiamiento de comercio o liquidación de valores, en lugar de un reemplazo total de la infraestructura existente. 
Este patrón de adopción parcial y heterogénea entre instituciones es consistente con la predicción del modelo dual de que la
asignación óptima $w_i^{A*}$ es, en general, estrictamente interior
($0 < w_i^{A*} < w_i^*$) y no una migración completa hacia la red Alternativa.

\section{Un modelo de redes duales}

El modelo base de \citeA{jalan_incentive-aware_2024} considera una red de $N$
agentes que negocian contratos bilaterales bajo una función de utilidad
media-varianza. Cada agente $i$ posee un vector de creencias $\mu_i$ sobre los
retornos esperados, una matriz de covarianzas $\Sigma_i$ y un coeficiente de
aversión al riesgo $\gamma_i$. La utilidad de un agente $i$ toma la forma
\begin{equation}
    g_i(W, P) = w_i^\top (\mu_i - P e_i) - \gamma_i\, w_i^\top \Sigma_i\, w_i,
\end{equation}
sujeta a una restricción regulatoria capturada por una matriz de proyección
$\Psi_i$. La solución óptima admite la forma cerrada $w_i^* = Q_i(\mu_i - Pe_i)$,
y una asignación $(W,P)$ es \textit{pairwise stable} si ningún par de agentes
puede mejorar simultáneamente desviándose unilateralmente de ella.

\citeA{venegas_reyes_financial_2025} extiende este marco introduciendo dos
infraestructuras paralelas. Primero, una red Tradicional ($T$) con restricción
regulatoria $\Psi_i^T$. Segundo, una red Alternativa ($A$) con restricción $\Psi_i^A$.
Cada agente reparte su exposición total entre ambas, $w_i^* = w_i^T + w_i^A$, y
la utilidad incorpora un beneficio de transparencia $\tau \geq 0$ de la red
alternativa y un costo friccional cuadrático $\lambda_A \geq 0$ sobre $w_i^A$:
\begin{equation}
    g_i(W, P) = w_i^\top (\mu_i + \tau e^A_i - P e_i) - \gamma_i\, w_i^\top \Sigma_i\, w_i
    - \lambda_A \|w^A_i\|^2.
\end{equation}
El modelo dual es consistente con el modelo base de Jalan y Chakrabarti en tres
casos límite. Cuando $\lambda_A \to 0$ y $\tau = 0$, cuando $\Psi^A$ desactiva
todos los contratos y finalmente cuando $\Psi^T = \Psi^A$.

\subsection{Resultado central}

Sobre una red de tres agentes, Banco Nacional ($N$), Banco Local ($L$) y Empresa
Local ($F$), donde la regulación prohíbe el contrato directo $N$--$F$,
\citeA{venegas_reyes_financial_2025} muestra que la sola competencia entre
intermediarios habilitados (incorporando una cuarta firma que compite con $L$)
recupera solo parcialmente el bienestar perdido por la restricción. Esta deja un
\textit{gap} de $27{,}9\%$ respecto del benchmark sin restricción. En cambio,
habilitar una red alternativa $A$ que sí permite el contrato $N$--$F$ logra una
recuperación completa del bienestar. La utilidad agregada del sistema iguala
exactamente al benchmark sin restricción, con el contrato $N$--$F$ transando casi
enteramente a través de $A$. Este resultado establece la red alternativa como el mecanismo estructural
que recupera eficiencia.

\subsection{Identificación y estimación desde datos}\label{c25:identificacion}

El objetivo específico 1 de este trabajo (Capítulo~\ref{c1}) es evaluar la
factibilidad de calibrar $\mu_i$, $\Sigma_i$ y $\gamma_i$ con datos reales.
Venegas (2024) no aborda esta pregunta: trabaja con esos primitivos dados como
parámetros ilustrativos, nunca estimados. El propio \citeA{jalan_incentive-aware_2024},
en cambio, sí ofrece resultados formales sobre qué tan factible es estimar estos
primitivos desde datos y bajo qué condiciones: resultados que Venegas no retoma,
pero que son el respaldo teórico directo de la validación empírica de este trabajo.

\paragraph{Estimación de la covarianza desde retornos históricos}
Un primer enfoque estima $\Sigma_i$ a partir del historial de retornos observado
por cada agente. Si el agente $i$ observa $m$ muestras independientes de los
retornos de contratos unitarios con los $n$ agentes de la red, recogidas en una
matriz $X_i \in \mathbb{R}^{n \times m}$, la covarianza muestral $\hat\Sigma_i$
converge en probabilidad a la covarianza poblacional $\Sigma$, pero con muestras
finitas cada agente puede terminar con una $\hat\Sigma_i$ distinta, lo que
degrada el factor de amortiguación seguro $\eta^*$ del Teorema~4 (arriba) y con
ello la velocidad de convergencia del Algoritmo~1. \citeA{jalan_incentive-aware_2024}
acotan qué tan severa es esa degradación: bajo condiciones de regularidad estándar
sobre $\Sigma$ y $\Gamma$ (normas de $\Sigma$, $\Sigma^{-1}$, $\Gamma$ y
$\Gamma^{-1}$ acotadas independientemente de $n$, y sin aristas prohibidas), si
cada agente estima su covarianza con $m = \lceil n \log n \rceil$ muestras
independientes, el factor de amortiguación estimado $\hat\eta^*$ se mantiene
arbitrariamente cerca del valor verdadero $\eta^*$ con probabilidad que tiende a
uno exponencialmente en $n$. La cota admite tamaños de muestra heterogéneos entre
agentes ($m_1,\dots,m_n$): basta con que el menor de ellos satisfaga
$\min_i m_i \geq \lceil n \log n \rceil$, lo que permite a instituciones con
distinto historial disponible usar ventanas de estimación de largo distinto.

La lectura correcta de este resultado no es que el requisito de datos disminuya
al crecer $n$ (aunque $\lceil n \log n \rceil$ sí crece con $n$, moderadamente),
sino que crece muy poco comparado con lo que exigiría estimar
sin estructura una matriz de covarianza $n \times n$ completa (con $n(n+1)/2$
entradas libres): la cota mínima teórica conocida para ese problema general es de
orden $n$ muestras, de modo que $n \log n$ queda solo un factor logarítmico por
encima del óptimo. Para la red de $n=3$ agentes de este trabajo
(BID/CORFO--intermediaria--MiPyME, Capítulo~\ref{MCP:subastador}), esto se
traduce en apenas $4$ a $5$ periodos históricos por agente según la base del
logaritmo, un requisito de datos modesto y no las decenas o cientos que una
estimación sin estructura podría sugerir.

\paragraph{Inferir creencias a partir de la estructura de la red}
Un obstáculo anterior a la pregunta de cuántos datos bastan es de
identificabilidad. \citeA{jalan_incentive-aware_2024} muestran que, dado un
único corte transversal de la red observada $W$ en su punto estable,
$(M,\Gamma,\Sigma)$ no pueden recuperarse de forma única, ni siquiera fijando
una escala (p.ej. $\operatorname{tr}(\Gamma)=\operatorname{tr}(\Sigma)=1$): para
cualquier elección válida de $\Gamma$ y $\Sigma$ existe una $M$ correspondiente
consistente con la misma $W$. Observar varias redes generadas con la misma
$\Sigma$ y $\Gamma$ pero distinta $M$ tampoco resuelve el problema. La
consecuencia práctica es directa: una sola observación de la red
BID/CORFO--intermediaria--MiPyME, por completa que sea, nunca bastaría para
calibrar el modelo. El dato requerido es necesariamente longitudinal, una
secuencia $W(t)$ observada en $T$ periodos distintos, no un corte transversal.

Para hacer tratable la inferencia sobre esa secuencia, \citeA{jalan_incentive-aware_2024}
introducen el Supuesto~1: (a) $\Gamma(t)=I$ y $\Sigma_i(t)=\Sigma$ para todo
$t\in[T]$, es decir, aversión al riesgo y covarianza homogéneas y constantes en el
periodo; (b) cada entrada $M_{ij}(t)$ evoluciona de forma independiente según
un movimiento browniano con los mismos parámetros para todos los pares
$(i,j)$; y (c) $\operatorname{tr}(\Sigma)=1$ como normalización de escala. El
supuesto (a) se motiva en la literatura de teoría de portafolios, donde los
errores de estimación de la media suelen dominar sobre los de la covarianza, de
modo que fijar $\Gamma=I$ y una $\Sigma$ constante concentra el problema de
estimación en $M$, que es lo que efectivamente varía en el tiempo. Nótese que
$\Gamma=I$ es, convenientemente, tanto el supuesto de homogeneidad de Venegas
(2024) como la condición que \citeA{jalan_incentive-aware_2024} necesitan para
que la inferencia de $\Sigma$ sea tratable, lo cual es compatible con la decisión
de alcance de este trabajo (Capítulo~\ref{c1}) de no abordar heterogeneidad en
$\gamma_i$ por ahora.

\paragraph{Proposición 1: estimar $\Sigma$ desde incrementos de la red, no contrato por contrato}
Bajo el Supuesto~1, \citeA{jalan_incentive-aware_2024} muestran que el
estimador de máxima verosimilitud de $\Sigma$ no requiere observar los retornos
subyacentes de cada agente (la ruta del Teorema~5), sino que se recupera
enteramente a partir de cómo cambia la red observada $W(t)$ de un periodo al
siguiente. La Proposición~1 reduce el problema a un programa semidefinido
(SDP):
\begin{equation}
    \min_{\Sigma} \sum_{t=1}^{T-1} \left\| \Sigma\big(W(t+1)-W(t)\big) + \big(W(t+1)-W(t)\big)\Sigma \right\|_F^2
    \quad \text{sujeto a } \Sigma \succeq 0,\ \operatorname{tr}(\Sigma)=1.
\end{equation}
La razón por la que este resultado es más relevante que el Teorema~5 para la
aplicación a datos de BID/CORFO es que la variable de entrada del problema es
directamente la secuencia de incrementos $\Delta W(t) = W(t+1)-W(t)$ del propio
contrato observado, es decir, los montos de financiamiento entre BID/CORFO, la
intermediaria y las MiPyMEs año a año, no una serie de retornos internos de
cada institución que en la práctica no está disponible públicamente. En vez de
estimar $\Sigma_i$ contrato por contrato a partir de datos que las instituciones
normalmente no publican, la Proposición~1 solo exige la misma serie de flujos de
financiamiento $W(t)$ que ya es el objeto de interés de la validación empírica
de este trabajo. El Remark~2 del paper generaliza además la Proposición~1 a una
$\Sigma$ que varía en el tiempo, agregando un regularizador
$\nu\sum_j \|\Sigma_{(j+1)}-\Sigma_{(j)}\|$ al objetivo del SDP que preserva la
convexidad del problema, lo cual sería relevante si el perfil de riesgo del
programa cambió estructuralmente durante 2016--2024 (p.ej. tras 2020).

\paragraph{Una limitación no resuelta por el paper, relevante para redes reguladas}
Ni el Teorema~5 ni la Proposición~1 discuten qué ocurre cuando la red tiene
aristas prohibidas por regulación, precisamente el caso de este trabajo, donde
el contrato directo BID/CORFO--MiPyME está excluido por diseño. Un experimento
propio con datos sintéticos generados sobre esta estructura de 3 agentes
(arista 1--3 prohibida) sugiere que la Proposición~1, tal como está formulada,
pierde identificabilidad en la entrada de $\Sigma$ correspondiente al par sin
contrato: como $W_{13}(t)\equiv 0$ para todo $t$ por construcción regulatoria,
esa entrada nunca aporta información al objetivo del SDP, y el error de
recuperación de $\Sigma_{13}$ no decrece aunque se aumente el número de periodos
observados $T$. Esta es una limitación metodológica propia de este trabajo, no
señalada por \citeA{jalan_incentive-aware_2024}, cuyos resultados de
identificación se enuncian para redes sin restricciones de arista, y se retoma
como hallazgo de factibilidad en el Capítulo~\ref{MCP:subastador}.

\subsection{Oportunidades de mejora identificadas}

El modelo de \citeA{venegas_reyes_financial_2025} obtiene su resultado de bienestar
bajo parámetros ilustrativos para $\mu_i$, $\Sigma_i$, $\gamma_i$ elegidos para
sostener los experimentos numéricos, no estimados de datos reales. Además, trabaja bajo dos
supuestos de homogeneidad entre agentes, con $\gamma_i = 1$ para todo agente y
$\mu_i = \mu_j$, que el propio autor reconoce como limitaciones. De las oportunidades de mejora identificadas (calibración empírica de los primitivos del modelo, incorporación de dinámica y dependencia de trayectoria, e incorporación de riesgo sistémico y contagio), este trabajo retoma la primera: calibrar el modelo con datos reales, como se detalla en el Capítulo~\ref{MCP:subastador}.
